% PAPER_A_METHOD_CONTRACT_CURRENT.json paper_a_method_20260807_v19 SHA-256 60e60f01d762ee89d788979ad9ef243db1887f6ef81f4d5bb8ae2e6f29b5af38
\section{V1：共享容量条件下的 Support-aware Full 政策试验}

\subsection{试验臂}

\begin{verbatim}
Strong Global
vs
Support-aware Full
\end{verbatim}

主要使用 Manual。

V1 使用独立的 V1 gate：\texttt{CALIBRATION\_ENROLLMENT\_CLOSED}、
\texttt{ALL\_CALIBRATION\_WORKERS\_TERMINAL}、\texttt{C1\_EVIDENCE\_FROZEN}、
\texttt{C2\_B\_FROZEN}、\texttt{C2\_A\_RP\_CLOSED}、\texttt{FINAL\_POOLED\_PROFILE\_FROZEN}、
\texttt{STRONG\_GLOBAL\_FROZEN}、\texttt{FULL\_POLICY\_FROZEN}、
\texttt{VALIDATION\_ROSTER\_FROZEN} 和 \texttt{V1\_SAP\_FROZEN}。
这些 roles 都必须绑定对应冻结 artifact、SHA 和依赖关系；C2-B closeout 只提供审计证据，
不能直接替代 Full Policy。

\subsection{Block 级共享容量合同}

每个 block 前冻结：

\begin{verbatim}
availability_snapshot
worker_total_capacity
global_quota_per_worker
full_quota_per_worker
candidate_roster
offer_timeout
completion_timeout
max_offer_attempts
replacement_rule
invalid_submission_replacement_rule
candidate_exhaustion_rule
\end{verbatim}

\subsection{两套独立容量账本}

\begin{verbatim}
capacity_global[u]
capacity_full[u]
\end{verbatim}

主分析不允许跨臂借用。

两臂必须共用：

\begin{verbatim}
offer timeout
completion timeout
maximum offer attempts
decline/no-response replacement
invalid submission replacement
candidate exhaustion terminal rule
\end{verbatim}

唯一差异是推荐排序。

\subsection{调度证据}

记录：

\begin{verbatim}
block_id
policy_arm
candidate_set_at_decision
availability_snapshot_id
policy_recommended_worker
recommendation_rank
offered_worker
offer_sequence
accepted_worker
completed_worker
replacement_reason
timeout
capacity_before
capacity_remaining
candidate_exhausted
\end{verbatim}

区分：

\begin{verbatim}
recommended_not_offered
offered_declined
accepted_not_completed
completed_invalid
completed_valid
\end{verbatim}

\subsection{动态追加}

初始与追加规则在两臂完全相同。

可冻结候选：

\begin{verbatim}
k_initial = 2
formal_structure_min_k = 3
standard_cap = 4
exceptional_cap = 5
\end{verbatim}

\(k_{\mathrm{initial}}=2\) 只表示首次 offer 数，不代表可以在两条合法提交时
形成正式 resolved 输出。crowd structure 与 stopping 至少要求三名不同且合资格的
worker IDs；未达到该数目时继续追加，至 cap 仍不足则记为 unresolved。实际 cap 与
其他调度数值由 C1/C2 replay 与功效/预算模拟冻结。

\subsection{GT-blind 聚合}

medoid 选择必须只读取合法提交、
冻结几何相似度、兼容簇、结构有效性和冻结 tie-break；不得读取 GT、公开 GT 评价或
依赖政策分数来决定 GT-blind medoid。

对合法提交：

\begin{enumerate}
  \item 将有冻结证据的 external system failure 移入重跑或行政删失处置，不作为工人或政策质量证据；
  \item 将 worker-caused structural failure 记录为 worker failure event，不进入合法聚合候选；
  \item 使用冻结几何相似度寻找最大内部一致簇；
  \item 判断 largest cluster、second cluster、medoid margin 和结构有效性；
  \item resolved 时选择主簇 medoid；
  \item medoid tie 依次使用 geometry SHA、canonical annotation identity 和
  task-bound frozen seed；这些 tie-break 不读取 Global、Full、GT 或 worker quality；
  \item 稳定多峰不强行输出；
  \item 到上限仍不稳定则 unresolved。
\end{enumerate}

\subsection{运行中间状态与最终政策终态}

\begin{verbatim}
external_system_failure_pending_disposition
\end{verbatim}

最终政策终态仅为：

\begin{verbatim}
resolved
unresolved
severe_failure
\end{verbatim}

\subsubsection{resolved}

产生合法 GT-blind 输出 \(O_t^\pi\)。

\subsubsection{unresolved}

存在合法提交，但未形成稳定单一输出。

\subsubsection{severe\_failure}

到上限仍没有可交付合法输出。

\subsubsection{external\_system\_failure\_pending\_disposition}

外部事故确认后不直接成为可比较的政策终态：按第 3.5 节在盲态下同臂、同冻结版本、对称预留容量重跑一次；不满足任一条件则行政删失。重跑不得跨臂借容量、不得使用不同冻结版本，也不得由 V1 结果触发。

\subsection{V1 outcomes}

\subsubsection{非交付指标}

\[
D_t^\pi
=
I(\text{unresolved or severe failure}).
\]

\subsubsection{Severe failure}

\[
H_t^\pi
=
I(\text{severe failure}).
\]

\subsubsection{Resolved-only quality}

\[
Q_t^\pi=IoU(O_t^\pi,G_t).
\]

\subsubsection{交付调整质量}

\[
U_t^\pi
=
I(\text{resolved})\times IoU(O_t^\pi,G_t).
\]

unresolved/severe failure 的 \(U_t^\pi=0\)，表示没有交付正式布局，不表示真实几何 IoU 必然为 0。

candidate exhaustion、替补规则失败和容量耗尽属于 policy-caused failure：保留在原政策臂 ITT，\texttt{policy\_failure=1}；若任务终态为 unresolved 或 severe failure，则交付调整质量为 0。worker invalid submission 作为 worker failure event 保留；若按冻结替补后任务 resolved，不把最终任务交付调整质量改写为 0。external system failure 的成功重跑以重跑任务保留在原臂 ITT；无法合规重跑的行政删失不进入质量分母，但必须单列原始、重跑、删失和事故原因的两臂分布。

严格集合口径为：randomized set 包含全部随机化任务；primary modified-ITT（mITT）
排除预注册且臂盲确认的 external administrative censor；unresolved、severe failure
和 policy-caused failure 仍留在原臂 mITT，交付调整质量为 0。每臂同时报告
randomized、rerun、administrative censor、mITT、resolved、unresolved 与 severe 数量。

\subsection{检验层级与确认性 margin}

\begin{enumerate}
  \item severe failure 不劣；
  \item unresolved＋severe failure 不劣；
  \item 交付调整政策质量；
  \item 标注数量或成本。
\end{enumerate}

resolved-only output quality、\(k_{used}\)、active time、完成时间和
policy\(\times\)risk\_route interaction 为次级结果。severe failure margin、
unresolved+severe margin、delivery-adjusted quality margin 及成本判定尺度均来自结果
可见前冻结的 Statistical Analysis Plan 与 policy analysis manifest；四级层级和 margin
不得由 V1 结果回调。容量、offer、replacement、dynamic redundancy、GT-blind aggregation
和 ITT 状态机的完整字段合同见附录 F--G。

\subsection{事后审查}

事后人工审查只用于：

\begin{itemize}
  \item 解释 unresolved；
  \item 识别 GT 冲突；
  \item 形成反例；
  \item 判断 OOS。
\end{itemize}

不能用人工修订后的输出替换政策输出重新计算 V1 主效应。

也不得在查看 T1/V1 结果后把 worker 或 policy failure 重分类为 external system failure，或选择性重跑某一臂。事故证据、分类时间和重跑/删失决定是独立审计对象。

\paragraph{Contract v5 clarification.}
The online engine consumes only the frozen \texttt{policy\_candidate\_v2.global\_rank\_S\_G}; it does not recompute peer or LOO ordering. Every start validates the method-contract SHA, policy-manifest SHA, candidate-roster SHA, profile version, and Stage 3 closure. The primary hierarchy is severe-failure non-inferiority, unresolved-plus-severe non-inferiority, delivery-adjusted-quality superiority, then count/cost.
